\section{Proof certificates with linear factors and a coefficient budget}
\label{sec:coefficient-budget}

This section develops a different arithmetic source for proof existence.
The proved contribution is a certificate theorem and its literal local
polynomial.  The final elimination of the unbounded trace has not been
materialized or measured.  Thus the local counts below are not a claim of a
new complete Diophantine-variable record.

Source~68 of the repository\cite{source68} already replaces ordered proofs by nonnegative
Horn flow and encodes formula labels with the pairwise coprime numbers
$1+Mj$.  It requires a second bounded universal block ensuring that $M$ is
divisible by all integers up to a label bound.  We remove that requirement.
We use the linear factors $X+j$, respect the actual arity of each rule, and
carry a separate, small arithmetic budget bounding polynomial coefficients.
The budget is chosen before the evaluation point.  No factorial, primality
test, common-multiple test, or variable exponent occurs in the new source.
The later elimination of bounded universal quantification may still use such
graphs; their costs have not disappeared from a final compiler by assertion.

\newtheorem{budgetlemma}{Lemma}[section]
\newtheorem{budgettheorem}[budgetlemma]{Theorem}
\newtheorem{budgetcorollary}[budgetlemma]{Corollary}

\subsection{Exact evaluation detects a polynomial identity}

All formula labels are natural numbers, including zero.  A rule has one
conclusion and zero, one, or two premises.  An occurrence of zero as an
actual premise is a genuine obligation.  In the C222 application, zero is
discharged by the explicit truth axiom; it is not discarded as a free
coordinate in the following construction.

For a finite list of rule instances
$R_i=(P_i\longrightarrow c_i)$, $0\leq i<N$, put
\begin{align}
 B_i&=\prod_{p\in P_i}(1+p)+c_i,
 &T_0&=1,
 &T_{i+1}&=T_i B_i,                                      \label{eq:budget-update}\\
 C(X)&=\prod_{i<N}(X+c_i),
 &Q(X)&=\prod_{i<N}\prod_{p\in P_i}(X+p),
 &M&=T_N+1.                                               \label{eq:budget-polynomials}
\end{align}
Products count repeated premises with their multiplicities.  An empty
product is one.  If two numerical premise slots $p_i,q_i$ are used, pad
unused slots with zero.  Then
\[
 B_i=(1+p_i)(1+q_i)+c_i
\]
for a binary rule; the same expression specializes to $1+p_i+c_i$ for a
unary rule and to $1+c_i$ for an axiom.

\begin{budgetlemma}[Coefficient budget]
Every coefficient of $C(X)-(X+1)Q(X)$ has absolute value at most $T_N$.
Consequently,
\begin{equation}
 C(M)=(M+1)Q(M)
 \quad\Longleftrightarrow\quad C(X)=(X+1)Q(X)\text{ in }\mathbb Z[X].
 \label{eq:budget-evaluation}
\end{equation}
\end{budgetlemma}
\begin{proof}
Each $B_i$ is at least both $1+c_i$ and
$\prod_{p\in P_i}(1+p)$.  Thus
$C(1)\leq T_N$ and $Q(1)\leq T_N$.
Write $Q(X)=\sum_j q_jX^j$, with $q_j\geq0$ and missing coefficients zero.
The coefficient of $X^j$ in $(X+1)Q(X)$ is $q_j+q_{j-1}$.
These are distinct entries in the coefficient list, so their sum is at
most $Q(1)$, including at either endpoint.  Every coefficient of $C$ is
at most $C(1)$.  The difference of two nonnegative numbers each at most
$T_N$ has absolute value at most $T_N$, not $2T_N$.

It remains to use a simple injectivity argument.  If a nonzero integer
polynomial $H(X)=\sum h_jX^j$ has $|h_j|<M$, let $k$ be the least index
with $h_k\ne0$.  Dividing $H(M)=0$ by $M^k$ and reducing modulo $M$ gives
$M\mid h_k$, contradicting $0<|h_k|<M$.  Apply this to
$H=C-(X+1)Q$ and $M=T_N+1$.  The reverse implication is immediate.
\end{proof}

The factor two sometimes used in such a coefficient bound is unnecessary.
Also, the product
$\prod_i(1+c_i)(1+p_i)(1+q_i)$ is an unnecessarily large budget.
The factor $B_i$ in~\eqref{eq:budget-update} gives the required two
majorizations separately and lowers the budget update from degree four to
degree three in its displayed coordinates.

\begin{budgetlemma}[Linear-factor uniqueness]
The polynomial identity in~\eqref{eq:budget-evaluation} holds exactly when
\begin{equation}
 \sum_{i<N}\left(e_{c_i}-\sum_{p\in P_i}e_p\right)=e_1
 \quad\text{in the free abelian group on }\mathbb N.
 \label{eq:budget-flow}
\end{equation}
\end{budgetlemma}
\begin{proof}
The distinct monic linear factors $X+j$, $j\in\mathbb N$, have distinct
roots.  Their multiplicities in an identity of products are equal.
This can be proved without invoking a general factorization theorem:
cancel the common copies of $X+j$ from the two sides and evaluate at
$X=-j$; a remaining copy on only one side would make exactly that side
zero.  This applies also to $j=0$, whose factor is $X$.
The resulting multiplicity equations are precisely
\eqref{eq:budget-flow}.
\end{proof}

\subsection{Why cycles cannot manufacture a refutation}

Let $\mathcal R$ be any collection of single-conclusion rules with at most
two premises and labels in $\mathbb N$.  Let $D$ be its least set of
derivable labels, starting from the zero-premise rules.

\begin{budgetlemma}[Single-conclusion flow]
Label $1$ is derivable if and only if some finite list of instances from
$\mathcal R$, with repetitions permitted, satisfies
\eqref{eq:budget-flow}.
\end{budgetlemma}
\begin{proof}
Unroll a finite proof into a tree.  Every non-root node contributes once
as a conclusion and once as a premise of its parent, leaving only the
root vector $e_1$.

Conversely, for a proposed finite list, form the finite closure of just
the rules in that list.  If $1$ is outside this closure, sum the coordinates
of the alleged flow equation that correspond to labels outside the closure.
Whenever a conclusion is outside, at least one premise is outside, since
the set is closed under the listed rule.  The contribution of that rule
to this sum is therefore at most zero.  When its conclusion is inside,
its contribution is again at most zero.  The total on the left is at most
zero, whereas the right side contributes one.  This is impossible.
The finite closure also gives a constructive proof DAG, which can then be
unrolled to a finite tree.
\end{proof}

Repeated premises, self-loops, and disconnected cyclic components do not
invalidate the proof.  The single-conclusion condition is essential: a
rule $a\longrightarrow(a,1)$ would have net vector $e_1$ without a way to
derive either conclusion from no axioms.

\begin{budgettheorem}[Two-accumulator certificate]
Label $1$ is derivable in $\mathcal R$ if and only if there exist $N,M$,
a list of $N$ valid rule instances $P_i\longrightarrow c_i$, and natural
sequences $S_0,\ldots,S_N$, $T_0,\ldots,T_N$ such that
\begin{align}
 S_0&=M+1, & S_N&=1, & T_0&=1, & M&=T_N+1,                \label{eq:budget-ends}\\
 S_{i+1}(M+c_i)&=S_i\prod_{p\in P_i}(M+p),
 &&&& (0\leq i<N),\nonumber\\
 T_{i+1}&=T_i\left(\prod_{p\in P_i}(1+p)+c_i\right),
 &&&& (0\leq i<N).                                       \label{eq:budget-transitions}
\end{align}
No external bound on the labels is required.
\end{budgettheorem}
\begin{proof}
For soundness, the budget recurrence implies $T_i\geq1$ and hence $M\geq2$.
Every factor $M+j$ is positive.  Starting from $S_0=M+1>0$, the first
recurrence shows inductively that every $S_i$ is positive.  Multiply all
of these recurrences and cancel the nonzero internal accumulators.
The endpoints give
$C(M)=(M+1)Q(M)$.
The coefficient-budget lemma, linear-factor uniqueness, and the flow
lemma now prove derivability of $1$.

For completeness, start with a finite proof tree.  List its nodes in a
top-down order in which the next expanded node is on the current frontier.
First compute the entire sequence $T_i$ from this list of labels; set
$M=T_N+1$ only afterwards.  Let $S_i$ be the product of $M+j$ over the
open frontier at stage $i$.  Initially the frontier is just label $1$;
at the end it is empty.  Expanding a node labelled $c_i$ removes its
factor $M+c_i$ and inserts the factors for its actual premises.  This
gives both recurrences and all four endpoints.  A zero-labelled leaf is
expanded with its actual zero-premise truth axiom when that axiom is used.
There is no circular choice of a scale or upper bound.
\end{proof}

The empty list cannot certify target $1$: it would require simultaneously
$M=2$ and $S_0=S_N=1=M+1$.  The argument is entirely finitary and uses
ordinary induction, integer arithmetic, and a finite closure construction.
It can be carried out within ZFC, without a standard-model soundness
assumption or an assumption that ZFC is consistent.

\subsection{The local arithmetic rows are quadratic after branch sharing}

Write the current and next accumulator values as $S,S',T,T'$, and use
$p,q,c$ for the two premise slots and conclusion.  The local flow and
budget constraints are as follows:
\begin{center}
\begin{tabular}{lll}
\hline
Rule arity & Flow equation & Budget equation\\
\hline
Zero & $S'(M+c)=S$ & $T'=T(1+c)$\\
One & $S'(M+c)=S(M+p)$ & $T'=T(1+p+c)$\\
Two & $S'(M+c)=S(M+p)(M+q)$
    & $T'=T(1+p+q+pq+c)$\\
\hline
\end{tabular}
\end{center}
For the explicit truth axiom $c=0$, the first row specializes further to
$S'M=S$ and $T'=T$.

Only the binary row needs additional local arithmetic wires.  Introduce
\[
 u=S(M+p),\qquad w=pq.
\]
Then its four residuals are
\begin{equation}
 u-S(M+p),\quad S'(M+c)-u(M+q),\quad w-pq,\quad
 T'-T(1+p+q+w+c).                                        \label{eq:budget-four-rows}
\end{equation}
They are all quadratic.  The witnesses $u,w$ are uniquely determined and
natural whenever the displayed input coordinates are natural.

The supplied local JSON augments the literal branch rows of
source~68's \texttt{horn\_shared\_source.json}.  Its binary modus-ponens
branch originally uses no scratch wires, so $u,w$ occupy its existing
shared slots \texttt{q0}, \texttt{q1}.  No active branch loses a wire or
shares two distinct values in one slot.  The most demanding inherited
branch still uses thirteen slots.  The generalization branch retains its
own original \texttt{q0}; it needs no new arithmetic wire.  The ten
zero-premise branches, including truth, need no new arithmetic wire either.

For each branch, add the displayed residual squares to that branch's
existing sum of squares.  Multiply the twelve resulting nonnegative
branch polynomials.  Over the naturals it vanishes exactly when one valid
branch and its corresponding accumulator update hold.  The construction
has twenty local existential coordinates:
\[
 p,q,c,\quad a,b,c_{\rm schema},d,\quad q_0,\ldots,q_{12},
\]
where the JSON uses \texttt{v} for the conclusion to avoid the displayed
name collision.  Its five exposed coordinates are $S,S',T,T',M$.
The truth factor rises from degree two to degree four, the inherited
bundle factor stays at degree eight, and all other factors have degree
four.  The local polynomial therefore has exact degree
$11\cdot4+8=52$.  This follows both from exact sparse expansion of every
factor and from nonvanishing of products over $\mathbb Z$.

The following comparison uses exactly the same inherited Horn rows and
the same unflattened branch-product representation.  It does not compare
the size of two final Diophantine polynomials.
\begin{center}
\begin{tabular}{lrr}
\hline
Source-interface quantity & Source 68 & This construction\\
\hline
Local existential coordinates & 23 & 20\\
Exposed local coordinates & 4 & 5\\
Local polynomial degree & 50 & 52\\
Outer certificate coordinates & 5 & 4\\
Bounded universal families & 2 & 1\\
Beta values needed at one step & 2 & 4\\
\hline
\end{tabular}
\end{center}
The extra beta values are a real cost: they encode the second accumulator.
Their arithmetic graph witnesses must be counted before asserting a net
improvement in a compiled polynomial.

\subsection{A completely specified trace interface}

Use four outer natural coordinates $N,x,y,M$ and
\[
 \beta(x,y,j)=x\bmod(1+(j+1)y),\qquad
 S_i=\beta(x,y,2i),\qquad T_i=\beta(x,y,2i+1).
\]
Impose the four endpoint equations~\eqref{eq:budget-ends}.  For every
$i<N$, require a zero of the supplied local polynomial with exposed
coordinates $(S_i,S_{i+1},T_i,T_{i+1},M)$.  This is one bounded universal
family; there is no additional family expressing divisibility of a scale
by all labels.

Every finite pair of accumulator sequences has such a common beta code.
For example, choose $y$ larger than all sequence entries and divisible by
$1,2,\ldots,2N+1$.  The moduli $1+(j+1)y$, $0\leq j\leq2N+1$, are
pairwise coprime: a common divisor is coprime to $y$, divides the index
difference, and that difference divides $y$.  The Chinese remainder
theorem supplies $x$, and the choice of $y$ makes the desired residues
their least nonnegative remainders.  This choice is made after $M$ and
the two traces; it creates no circular scale dependency.

For the calibrated C222 rule relation, the resulting finite-certificate
predicate is equivalent to existence of a C222 refutation of label $1$.
Composing with that calculus's inherited syntactic calibration gives the
same consistency target.  This route checks rule instances directly; its
source does not simulate K215 and does not use a Jones finalizer.

The remaining work is precise.  Substitute literal remainder graphs for
the four beta values, apply a certified bounded-universal elimination to
the unflattened local source, quantify the four outer coordinates, and
only then choose a final degree-reduction or variable-compression
strategy.  The full ordinary polynomial, its complete variable count,
and its coefficient payload have not been produced here.  The theorem
above is not a substitute for those artifacts.

\subsection{An optional normalized variant}

If one wishes to retain the earlier free-zero convention, use weights
$1+jX$ and pad every rule to two premises.  Keep
$T'=T((1+p)(1+q)+c)$ and $M=T_N+1$.  The same coefficient proof applies
to products of $1+jX$, and factor uniqueness recovers all positive-label
multiplicities; the factor for zero is one.  The soundness theorem then
uses zero as a designated derivable label.

The flow equation is
\[
 S'(1+Mc)=S(1+Mp)(1+Mq).
\]
It has a useful exact quadratic lift with four natural auxiliary values:
\[
 w=pq,\quad u=MS,\quad u'=MS',\quad v=Mw.
\]
The six quadratic residuals are
\begin{gather*}
 w-pq,\quad u-MS,\quad u'-MS',\quad v-Mw,\\
 S'+cu'-S-u(p+q+v),\qquad T'-T(1+p+q+w+c).
\end{gather*}
Their sum of squares is quartic, and its zeros are exactly the two desired
updates together with the uniquely determined natural auxiliaries.
Expanding the fifth residual gives
$S'(1+Mc)-S(1+Mp)(1+Mq)$ exactly.  This optional construction explains
how to reuse the previous padded Horn-flow interface, but the unnormalized
linear-factor source is cheaper locally because it respects rule arity.

\subsection{Checks and a necessary negative control}

The standalone checker reconstructs all twelve branch sum-of-squares
polynomials exactly from their literal residuals.  It checks 178 exact
natural branch fixtures against independently written C222 constructor
formulas, including both signs of the guarded bundle.  It exhausts all
60,880 lists of up to three zero-, one-, or two-premise rules on labels
$\{0,1,2\}$, checking the coefficient bound, numerical identity,
polynomial identity, flow equation, and derived-target implication.
Four toy proof-tree fixtures cover repeated premises and explicit truth
leaves; they are finite artificial Horn systems, not ZFC refutations.
The optional normalized lift has its own exact arithmetic checks.

The budget endpoint is essential.  Take $M=11$ and the two rules
$10\longrightarrow3$ and the axiom $7$.
The target $1$ is underivable, but
\[
 (M+3)(M+7)=(M+1)(M+10),\qquad 14\cdot18=12\cdot21.
\]
Indeed $S=(12,18,1)$ satisfies both local flow updates, and $M$ is larger
than every label.  The budget is $T=(1,14,112)$, so the required endpoint
would be $M=113$, not $11$.  The complete checker rejects the fake.
Thus an unspecified assertion that the evaluation point is ``large''
cannot replace the coefficient proof.

These finite checks are diagnostic evidence for the implementation.
The general mathematical justification is the preceding proof, and no
kernel-checked formalization of this new source is claimed.
